% DRAFT: the fresh Cartier export and its independent replay are still pending.
% No statement in this file is entered in the accepted ledger until the final
% certificate manifest marks all required stages complete.
\subsection{Excluding a point of order $49$}
\begin{proposition}\label{prop:decic-no49}
An elliptic curve over a number field of degree $10$ has no point of order $49$.
\end{proposition}
\begin{proof}
Put $X=X_1(49)$ and $J=J_1(49)$. By
\cite[Theorems~3.1 and 4.13, and Table~2]{DEHMZ},
\[
 J(\mathbb Q)\simeq C_7\times C_{52367710906884085342}.
\]
In particular $J(\mathbb Q)[2^\infty]\simeq C_2$.
The map $\pi:X\to X_0(49)$ has degree $21$ and diamond group
$H\simeq C_{21}$. The latter acts trivially on the rational
$2$-primary subgroup, so
$\pi^*\pi_*t=\sum_{h\in H}ht=21t$ makes $\pi_*$ injective there.
Similarly $\pi_*\pi^*=[21]$ makes pullback injective on
$J_0(49)(\mathbb Q)[2^\infty]$, whose order is therefore at most two.
The model
\[
 y^2+xy=x^3-x^2-2x-1
\]
for $J_0(49)$ has discriminant $-343$ and a point of order two
$(2,-1)$, reducing to the nonzero point $(0,1)$ modulo $2$.
Thus reduction is injective on the rational $2$-primary subgroup
of $J_0(49)$, hence also on that of $J$. It is always injective on
odd-order torsion. Therefore
\[
 J(\mathbb Q)\longrightarrow J(\mathbb F_2)
\]
is injective.

Suppose that $D$ is a rational effective noncuspidal divisor of
degree $10$ on $X$. Exhaustive enumeration over $\mathbb F_{2^f}$,
$f\leq10$, shows that the only possible noncuspidal closed-point
degrees are $9$ and $10$. There are respectively $42$ and $126$
such points, forming two and six diamond orbits. For every orbit
representative $P$, the degree-zero divisor
\[
 A_3P=(T_3-1-3\langle3\rangle)P
\]
is nonprincipal, by exact Riemann--Roch computations on the full
genus-$69$ curve. Independent enumeration of all curves, twists
and level structures verifies the completeness of these orbits.
The four native $3$-isogenies at each of the eight representatives
also reproduce the finite correspondence, including all closed
point degrees and norm multiplicities.

The operator $A_3$ kills $J(\mathbb Q)$, by Eichler--Shimura at $3$
and injectivity of torsion reduction there. At a rational cusp $c$
we have the literal identity $T_3c=c+3\langle3\rangle c$.
It follows from the Tate parametrization at one rational cusp and
then from diamond equivariance at all $21$ rational cusps.
If the reduction of $D$ has noncuspidal support, it is therefore
one degree-$10$ point or a degree-$9$ point plus a rational cusp.
In either case the computed nonprincipality of $A_3P$ is a
contradiction.

The cusp degrees over $\mathbb Q$ and $\mathbb F_2$ are
\[
 1^{21},6^3,21\qquad\text{and}\qquad1^{21},3^6,21.
\]
If the reduction of $D$ is supported on rational cusps, let $C$
be its rational cuspidal lift. Injectivity of specialization gives
$D\sim C$. Since
\[
 \operatorname{gon}_{\mathbb C}X
 \geq \frac7{800}[\operatorname{PSL}_2(\mathbb Z):\overline{\Gamma_1(49)}]
 =\frac7{800}\,1176=\frac{1029}{100}>10,
\]
we obtain $D=C$, a contradiction.

It remains to exclude cuspidal reductions containing a cubic cusp.
Put $U=(\mathcal C-1)H^0(X_{\mathbb F_2},\Omega^1)$, where
$\mathcal C$ is Cartier. Its dimension is $68$. We verify that
restriction from $U$ has rank $10$ at every effective cuspidal
divisor of degree $10$ containing a cubic cusp, as follows.

First use the existing common-basis jets at eighteen rational
cusps, through multiplicity six, and at all six cubic cusps,
through multiplicity three. The complete degree-$10$ replay checks
\[
 6\left(\binom{24}{7}-18\right)
 +\binom72\binom{21}{4}+\binom83\,18
 =2203209
\]
divisors and finds rank $10$ in every case. This covers every
mixed cuspidal divisor having at most six rational support points
and rational multiplicities at most six: the diamond group acts
regularly on the $21$ rational cusps, and at most $6\cdot3=18$
of its $21$ elements send some support point to one of the three
omitted cusps. Cartier commutes with diamonds.

If there are seven rational support points, the divisor consists
of seven distinct rational cusps and one cubic cusp. Fresh
first-order jets at all $21$ rational and six cubic cusps check
all $6\binom{21}{7}=697680$ such divisors. Finally, a rational
multiplicity of seven forces the divisor to be $7c+Q$, with $Q$
a cubic cusp. A diamond fixes the choice of $c$, leaving six
checks, using fresh jets through order seven at that cusp.
All these ranks are $10$. A separate calculation with Magma's
built-in local completion verifies the restriction kernels through
all seven rational orders and at a cubic cusp, including their
combined degree-$10$ kernel. Binary linear algebra independently
checks every exported projection and every stated rank.

Lemma~\ref{lem:decic49-Cartier}, applied with the full Jacobian as
target, now excludes the remaining cuspidal reductions. Each
closed cubic cusp is lifted with its complete Frobenius orbit,
so the entire cuspidal lift is a $\mathbb Q_2$-rational divisor.
This completes the proof.
\end{proof}
